We investigate the extent to which contemporary Large Language Models (LLMs) can engage in \emph{exploration}, a core capability in reinforcement learning and decision making. We focus on native performance of existing LLMs, without training interventions. We deploy LLMs as agents in simple \emph{multi-armed bandit} environments, specifying the environment description and interaction history entirely \emph{in-context}, \ie within the LLM prompt. We experiment with \gptthree, \gptfour, and \llama, using a variety of prompt designs, and find that the models do not robustly engage in exploration without substantial interventions: i) Across all of our experiments, only one configuration resulted in satisfactory exploratory behavior: \gptfour with chain-of-thought reasoning and an externally summarized interaction history, presented as sufficient statistics; ii) All other configurations did not result in robust exploratory behavior, including those with chain-of-thought reasoning but unsummarized history.
Although these findings can be interpreted positively, they suggest that external summarization---which may not be possible in more complex settings---is important for obtaining desirable behavior from LLM agents.
We conclude that non-trivial algorithmic interventions, such as fine-tuning or dataset curation, may be required to empower LLM-based decision making agents in complex settings.




